\documentclass[11pt]{article}
\usepackage[utf8]{inputenc}
\usepackage[T1]{fontenc}
\usepackage[italian]{babel}
\usepackage{geometry}
\geometry{a4paper, margin=1in}
\usepackage{xcolor}
\usepackage{listings}
\usepackage{hyperref}

\usepackage{tcolorbox}
\tcbuselibrary{breakable}
\begin{document}

\begin{tcolorbox}[colback=blue!5!white,colframe=blue!75!black,title=\textbf{DOMANDA}, breakable]
sommatrice UDP e influenze reciproche\\
10)
\begin{itemize}
    \item Invocare il server della sommatrice con due client diversi (tutti e tre possono anche essere sulla stessa macchina ovviamente su finestre terminali diverse), ad es.
    \item Che somma leggo da ciascun client? E' la somma che ciascun client da solo si aspetterebbe?
\end{itemize}

Client A:
\begin{itemize}
    \item Digitare "2345" + INVIO
    \item Digitare "5187" + INVIO
    \item Digitare "2" + INVIO
    \item Digitare "1" + INVIO
    \item "0"
\end{itemize}

Client B:
\begin{itemize}
    \item Digitare "2" + INVIO
    \item Digitare "8" + INVIO
    \item "0"
\end{itemize}
\end{tcolorbox}

\begin{tcolorbox}[colback=green!5!white,colframe=green!75!black,title=\textbf{RISPOSTA}, breakable]
In questo scenario, il server si comporterà in modo anomalo dal punto di vista applicativo.

Ecco le risposte alle due domande:

\begin{itemize}
    \item \textbf{Che somma leggo da ciascun client?} Il client che per primo invia il valore "0" leggerà una somma "mista" composta dai numeri inviati da entrambi i client fino a quel momento. Supponendo che i pacchetti arrivino in modo intervallato, il server sommerà tutti i valori insieme (ad esempio 2345 + 5187 + 2 + 1 + 2 + 8 = 7545). Il primo client a inviare lo "0" riceverà come risposta 7545 (o un valore intermedio a seconda di quanti pacchetti di B sono arrivati prima del suo "0"). Il secondo client (che manda lo "0" per ultimo) non riceverà nulla se il server termina l'esecuzione, o riceverà una somma azzerata/parziale se il server ha un ciclo esterno che lo fa ricominciare.
    
    \item \textbf{È la somma che ciascun client da solo si aspetterebbe?} Assolutamente no. Il Client A si aspetterebbe 7535 e il Client B si aspetterebbe 10. Invece, i dati si sono mischiati.
\end{itemize}

\textbf{Perché succede questo?} A differenza di \texttt{TCP} (che crea un "tubo" dedicato e isolato per ogni connessione/client), \texttt{UDP} non ha il concetto di connessione. Il server ha una sola coda di ricezione per la porta \texttt{UDP} e una singola variabile somma. Quando esegue \texttt{UDPReceive}, legge semplicemente il prossimo pacchetto in coda, senza distinguere se proviene dal Client A o dal Client B, sommandolo all'accumulatore globale. Inoltre, quando il server riceve lo "0" da uno dei due, il ciclo si interrompe e la funzione \texttt{UDPSend} risponde all'indirizzo \texttt{IP} e alla porta dell'ultimo pacchetto ricevuto (cioè quello del client che ha mandato lo "0"), inviando il "calderone" di numeri mischiati.

Questo esercizio evidenzia che \texttt{UDP}, non gestendo nativamente le sessioni, richiede allo sviluppatore di identificare esplicitamente da quale client proviene ogni messaggio e di gestire le variabili di stato (come la "somma") separatamente per ogni utente.
\end{tcolorbox}

\end{document}
